\section{Implementazione}
\label{sec:implementazione}

\subsection{Tecnologie e struttura}

Il programma è scritto in \textbf{Python 3.11+} con interfaccia grafica \textbf{PyQt6}. Le entità sono \code{dataclass}, le interfacce sono classi astratte (\code{abc}), gli importi usano \code{Decimal} per evitare errori di arrotondamento e i dati sono salvati in un file JSON. I test usano \code{unittest} (PyUnit); i diagrammi sono stati fatti con Enterprise Architect.

Il codice è in \code{src/ticketseller}, con una cartella per ogni parte (\code{views}, \code{controllers}, \code{models}, \code{dao}, \code{patterns}, \code{external_services}). Il file \code{app.py} crea gli oggetti, li passa al controller e apre la finestra richiesta:

\par\noindent\begin{minipage}{\linewidth}
\begin{lstlisting}
python -m ticketseller.app kiosk     # postazione del cliente
python -m ticketseller.app clerk     # postazione del cassiere
\end{lstlisting}
\end{minipage}\par

Negli estratti seguenti sono stati tolti commenti, docstring e alcune annotazioni di tipo.

\subsection{Controller}

Il controller riceve dall'esterno tutto quello che usa (DAO, servizio di pagamento, strategie di prezzo): così nei test si possono passare oggetti di prova.

\par\noindent\begin{minipage}{\linewidth}
\begin{lstlisting}
class TicketSeller:
    reservation_timeout = timedelta(minutes=15)

    def __init__(self, event_dao, ticket_dao, payment_service, pricing_strategy,
                 subscription_strategy=None, subscription_dao=None):
        self.event_dao = event_dao
        self.ticket_dao = ticket_dao
        self.payment_service = payment_service
        self.pricing_strategy = pricing_strategy
        self.subscription_strategy = subscription_strategy
        self.subscription_dao = subscription_dao
        self._reservation_lock = RLock()
\end{lstlisting}
\end{minipage}\par

\subsection{Acquisto}

\code{purchase_access_title} esegue l'acquisto (UC3): blocca i posti, calcola il totale, chiede il pagamento e poi conferma la vendita o libera i posti. Se i dati della carta non sono validi il pagamento solleva un errore: prima di passarlo alla finestra, i posti vengono liberati.

\par\noindent\begin{minipage}{\linewidth}
\begin{lstlisting}
def purchase_access_title(self, event_id, seat_ids, customer, payment_data, subscription=None):
    tickets = self.reserve_seats(event_id, seat_ids, customer)
    try:
        total = self.calculate_total(tickets, subscription)
        result = self.payment_service.process_payment(payment_data, total)
    except ValueError:
        self.cancel_reservation(tickets)
        raise
    if result.approved:
        self.confirm_sale(tickets, customer)
    else:
        self.cancel_reservation(tickets)
    return result
\end{lstlisting}
\end{minipage}\par

In \code{reserve_seats} il controllo dei posti e il blocco avvengono sotto un lock, così due richieste contemporanee non possono prendere lo stesso posto (RNF3):

\par\noindent\begin{minipage}{\linewidth}
\begin{lstlisting}
def reserve_seats(self, event_id, seat_ids, customer=None):
    with self._reservation_lock:
        self.release_expired_reservations(event_id)
        if len({str(seat_id) for seat_id in seat_ids}) != len(seat_ids):
            raise ValueError("Duplicate seats cannot be selected.")
        tickets = self.ticket_dao.find_tickets_by_seats(event_id, seat_ids)
        if not tickets:
            raise ValueError("At least one seat must be selected.")
        if any(not ticket.is_available() for ticket in tickets):
            raise ValueError("One or more selected seats are not available.")
        for ticket in tickets:
            ticket.reserve(customer)
        return self.ticket_dao.save_tickets(tickets)
\end{lstlisting}
\end{minipage}\par

\subsection{Pattern State}

Il biglietto non decide da solo cosa fare: passa la richiesta al suo stato corrente, che controlla se l'operazione è permessa e cambia stato.

\par\noindent\begin{minipage}{\linewidth}
\begin{lstlisting}
class Ticket(Voucher):
    def reserve(self, customer=None):
        self.state.reserve(self, customer)

    def cancel(self):
        self.state.cancel(self)

    def confirm_sale(self, customer):
        self.state.confirm_sale(self, customer)


class LockedState(TicketState):
    def reserve(self, ticket, customer=None):
        raise ValueError("Ticket is already locked.")

    def cancel(self, ticket):
        ticket.set_state(AvailableState())
        ticket.locked_at = None
        ticket.owner = None

    def confirm_sale(self, ticket, customer):
        if ticket.owner is not None and ticket.owner != customer:
            raise ValueError("Ticket is locked by another customer.")
        ticket.set_state(SoldState())
        ticket.locked_at = None
        ticket.owner = customer
\end{lstlisting}
\end{minipage}\par

\subsection{Pattern Strategy}

Ogni politica di prezzo è una classe con il metodo \code{calculate_total}. Il controller usa la strategia per gli abbonamenti quando il cliente ne indica uno, altrimenti quella normale.

\par\noindent\begin{minipage}{\linewidth}
\begin{lstlisting}
def calculate_total(self, tickets, subscription=None):
    strategy = self.pricing_strategy
    if subscription is not None and self.subscription_strategy is not None:
        strategy = self.subscription_strategy
    return strategy.calculate_total(tickets, subscription)


class SubscriptionPricing(PricingStrategy):
    def calculate_total(self, tickets, subscription=None):
        if subscription is None:
            raise ValueError("Subscription pricing requires a subscription.")
        discount_rate = subscription.discount_rate if subscription.discount_rate is not None else Decimal("1.00")
        total = Decimal("0.00")
        for ticket in tickets:
            if ticket.event_id in subscription.covered_event_ids:
                total += ticket.base_price * (Decimal("1.00") - discount_rate)
            else:
                total += ticket.base_price
        return quantize_money(total)
\end{lstlisting}
\end{minipage}\par

\subsection{Persistenza e rilascio dei posti}

\code{TicketDB} tiene i dati in memoria e a ogni modifica riscrive il file \code{database.json}; all'avvio lo rilegge. Il resto del programma non sa che i dati stanno in un file, perché usa solo le interfacce DAO: per passare a un database basterebbe un'altra classe che le realizza.

Il kiosk avvia un timer che ogni minuto libera i posti bloccati da più di 15 minuti (attore Tempo, UC6):

\par\noindent\begin{minipage}{\linewidth}
\begin{lstlisting}
self.release_timer = QTimer(self)
self.release_timer.timeout.connect(self._release_expired_reservations)
self.release_timer.start(60_000)
\end{lstlisting}
\end{minipage}\par
